\documentclass[10pt,a4paper]{amsart}
\usepackage[margin=19mm]{geometry}
\usepackage{amsmath,amssymb,mathtools}
\usepackage[scr=rsfs,cal=euler]{mathalfa}
\usepackage{xfrac}
\usepackage[unicode,hidelinks,bookmarks=false]{hyperref}
\newcommand{\ie}{i.e.\ }
\newcommand{\cf}{cf.\ }
\newcommand{\resp}{respectively\ }
\newcommand{\Subsection}{\subsection}
\providecommand{\nobreakdash}{\nobreak\hbox{-}}
\setlength{\emergencystretch}{2em}
\multlinegap0pt
\usepackage{amssymb}
\newtheorem{proposition}{Proposition}
\newtheorem{corollary}{Corollary}
\newcommand{\Dbb}{\mathbb D}
\newcommand{\bd}{{B_d}}
\newcommand{\cph}{C_\varphi}
\newcommand{\dbr}{\mathcal{H}}
\newcommand{\meal}{\sigma_d}
\newcommand{\om}{\omega}
\newcommand{\ph}{\varphi}
\newcommand{\spd}{\partial B_d}
\newcommand{\za}{\zeta}
\title{Composition operators between de Branges-Rovnyak and Hardy spaces}
\author{Evgueni Doubtsov, Andrei V. Vasin}
\date{Source: arXiv:2606.22965v1 (CC BY 4.0)}
\begin{document}
\maketitle

\noindent\emph{Source and licence.} Composition operators between de Branges-Rovnyak and Hardy spaces. Version arXiv:2606.22965v1 (CC BY 4.0), distributed by arXiv under the Creative Commons Attribution 4.0 International licence, http://creativecommons.org/licenses/by/4.0/, as recorded in the arXiv licence metadata for this exact version and shown on the abstract page https://arxiv.org/abs/2606.22965v1. Full licence notice: \texttt{licenses/arxiv-2606.22965-CC-BY-4.0.txt}.\par\smallskip

\noindent\textit{Editorial scope.} Counted unit: the article's corollary c\_suff (source0.tex lines 528--537). The article's complete original proof of it is delivered with it (source0.tex lines 538--555, one \texttt{proof} environment of 18 lines). No mathematics is written, repaired or re-derived: the statement and the proof are the article's own lines, in the article's own notation, and no retained line is substituted or re-flowed.  Retained uncounted as the source-local context the proof needs: lines 369--371; lines 412--420.  Nothing was omitted from any retained span, so the package carries no omission record.  No retained line contains a citation, so the wrapper carries no bibliography and no bibliography entry.  No retained line contains a figure environment, an image-inclusion command or drawing code, so no image is delivered and the package declares no asset dependency.  Editorial formatting only, in addition: an AMS \texttt{amsart} preamble that restates the article's own declarations and macros used by the retained lines, namely the environments \texttt{proposition}, \texttt{corollary} on separate counters, together with the standard packages those lines need.  The retained environments print in this document as Proposition 1, Corollary 1 in order of appearance, whereas the article prints the same items with the numbering of its own full document.  The complete native source of the article as distributed in the arXiv e-print for this exact version is bundled unchanged as \texttt{sources/203393/source0.tex}.
\begin{equation}\label{e_om_p_A}
\om_p(z)\ge A \frac{1-|z|}{\left(1- |b(z)| \right)^{\frac{p}{q(p+1)}}}, \quad z \in \Dbb.
\end{equation}

\begin{proposition}\label{p_suff}
Let $\ph: \bd \to \Dbb$ be a holomorphic function, $\ph(0)=0$.
Let $b:\Dbb \to \Dbb$ be a holomorphic function such that $\dbr(b)$ is infinite dimensional.
Assume that there exists $p\in (1,2)$ such that
\begin{equation}\label{e_om2Neva0}
\om_p^{-2}(z) \int_{\spd} N_{\ph_\za} (z)\, d\meal(\za) \to 0\quad \textrm{as}\ |z|\to 1-.
\end{equation}
Then the operator $\cph: \dbr(b) \to H^2(\bd)$ is compact.
\end{proposition}

\begin{corollary}\label{c_suff}
Let $\ph: \bd \to \Dbb$ be a holomorphic function, $\ph(0)=0$.
Let $b:\Dbb \to \Dbb$ be a holomorphic function such that $\dbr(b)$ is infinite dimensional.
Assume that there exists $\gamma\in (0, \frac{1}{3})$ such that
\begin{equation}\label{e_gamNeva0}
\left(\frac{(1-|b(z)|)^\gamma}{1-|z|^2} \right)^2 \int_{\spd} N_{\ph_\za} (z)\, d\meal(\za) 
\to 0\quad \textrm{as}\ |z|\to 1-.
\end{equation}
Then the operator $\cph: \dbr(b) \to H^2(\bd)$ is compact.
\end{corollary}

\begin{proof}
Put $p = \frac{1+\gamma}{1-\gamma}$. Observe that $p\in (1,2)$.
Let $q$ denote the conjugate exponent of $p$.
Then
\[
\gamma = \frac{p-1}{p+1} = \frac{p}{q(p+1)}.
\]
By \eqref{e_om_p_A}, we have
\[
\om_p (z) \ge A\frac{1-|z|}{(1-|b(z)|)^\gamma}.
\]
Therefore, \eqref{e_gamNeva0} implies that
\[
\om_p^{-2}(z) \int_{\spd} N_{\ph_\za}(z)\, d\meal(\za) \to 0
\]
as $|z|\to 1$. Applying Proposition~\ref{p_suff}, we conclude that the operator
$\cph: \dbr(b) \to H^2(\bd)$ is compact.
\end{proof}

\end{document}
